\documentclass[twocolumn,prd,superscriptaddress,nofootinbib]{revtex4-2}

\usepackage{amsmath,amssymb}
\usepackage{graphicx}
\usepackage{hyperref}

\begin{document}

\title{Radial Dimensionality Theory II: White Dwarf Structure under Solar-Constrained Density-Dependent Dimensionality}
\author{Christopher K. Merrill}
\affiliation{Independent Researcher}

\author{Claude (Anthropic)}
\affiliation{Large Language Model, Computational Collaboration}

\author{Astra (Open AI)}
\affiliation{Large Language Model, Computational Collaboration}

\date{Consistency revision: November 12, 2026}

\begin{abstract}
We propagate the revised nominal solar parameter sets of Paper~I through
an audited $n=3/2$, $\mu_e=2$ white-dwarf polytropic prescription. The
source-precision solar-a score has a sampled amplitude envelope .00874717,
locally refined to .00875020; these are conditional model/proxy thresholds,
not calibrated observational confidence bounds. We evaluate 66 exponent/
boundary points, 12 central densities and five amplitude fractions within
each nominal slice, totaling 3960 models. The largest sampled exact
fixed-mass radius deficit is .954933\%, at comparison mass
$M_A=.848321M_\odot$, $\alpha=.5$, $A=.00673493$ and
$\rho_c=3\times10^6~{\rm g\,cm^{-3}}$. This high-density point lies beyond
the non-relativistic EOS applicability estimate and is a mathematical
polytrope comparison. The figure uses the actual comparison mass and
$\exp(\Delta\ln R)-1$, rather than a baseline mass label and log shift.
The sampled response decreases with amplitude but does not certify a
global continuum envelope or universal contraction. Earlier V4 numerical
results remain superseded; realistic stellar responses and measurement
systematics are needed for quantitative empirical white-dwarf constraints.

\end{abstract}

\maketitle

\section*{Version note: v6}
This version retains the published title. It propagates the revised solar parameter sets and corrects the fixed-mass comparison and exact percentage shifts. The results describe sampled polytropic white-dwarf models within the stated equation-of-state limits; they do not establish general predictions for realistic white dwarfs.


\noindent\textbf{Scope.}
$A_{\Delta Q=2.71}(\alpha)$ denotes Paper~I's nominal fixed-exponent endpoint,
with $\rho_0=150~{\rm g\,cm^{-3}}$ and the solar-a reference prescription.
Its sampled/local maxima .00874717/.00875020 have no calibrated confidence
coverage. The asymptote associated with the latter is approximately 2.97375;
finite-density dimensions require the full law. The same input parameters
specify this white-dwarf model without retuning, but do not establish the
accuracy of its EOS or modified response. Empirical tests of a specified
phenomenological model do not require external endorsement.

\section{Introduction}
%============================================================================

Paper~I of this series \cite{PaperI} formulated Radial Dimensionality
Theory (RDT) as a phenomenological law in which the effective spatial
dimensionality felt by the gravitational field decreases weakly with local
density,
\begin{equation}
\begin{gathered}
d_{\rm eff}(\rho)=3\,\Omega(\rho),\quad
\Omega(\rho)=1-A\,\frac{x}{1+x},\\
x=\left(\frac{\rho}{\rho_0}\right)^{\!\alpha}.
\end{gathered}
\label{eq:law}
\end{equation}
with $\rho_0 = 150~{\rm g\,cm^{-3}}$ fixed as a series convention, and
derived the corresponding exact modified Lane--Emden equation with the
modification confined to the gravitational divergence operator. Confronting
solar neutrino fluxes and a helioseismic proxy bound, Paper~I found only a small combined-score improvement over the null and
obtained a nominal sampled/local amplitude envelope
\begin{equation}
A_{\rm sampled}=.00874717,\quad A_{\rm local}\simeq.00875020,
\label{eq:solarbound}
\end{equation}
for sampled exponents $\alpha\in[.02,.5]$, without confidence coverage. The associated high-density asymptote is approximately 2.97375.

This paper asks what that bound permits in the next density regime:
white dwarfs, with central densities $\rho_c \sim 10^5$--$10^6~{\rm
g\,cm^{-3}}$, three to four orders of magnitude beyond the solar core.
Because the law of Eq.~(\ref{eq:law}) is a function of density alone, the
solar constraint propagates to white dwarfs with \emph{no retuning}: each
chosen parameter point specifies a white-dwarf calculation under the
adopted polytropic assumptions. We sample each nominal slice at five amplitude fractions, including its endpoint; a global continuum envelope still requires a separate domain check.

This version supersedes an earlier analysis (V4) of white dwarfs in this
series. As documented in Sec.~\ref{sec:v4}, the V4 numerical
implementation was invalid, and its modification ansatz differed from the
operator form adopted by the series; neither its quantitative predictions
(percent-level radius \emph{increases}) nor its Chandrasekhar-mass shifts
survive. The retained operator calculation gives radius \emph{deficits}
near one percent at fixed mass in the sampled polytropes; the numerical
audit below supports those samples, not a global observational bound.

%============================================================================
\section{Framework}
%============================================================================

\subsection{Modified structure equation}

Following Paper~I, physical space remains three-dimensional for volume
elements, mass integrals, and microphysics; only the gravitational
divergence operator carries $d_{\rm eff}$, with reference scale
$r_\star = a$ (the Lane--Emden length) so that the model is
unit-invariant. This fixes a physical prescription: changing $r_\star$
at fixed stellar structure would change the variable-$\delta$ operator.
The length $a$ is spatially constant for each solution but changes with
$\rho_c$ across the family. Density here is Newtonian matter mass per
ordinary three-dimensional volume. For a polytrope $P = K\rho^{1+1/n}$, $\rho = \rho_c
\theta^n$, $\xi = r/a$, the exact structure equation is
\begin{equation}
\theta'' + \left[\frac{2+\delta(\xi)}{\xi}
+ \frac{d\delta}{d\xi}\ln\xi\right]\theta' + \theta^{n} = 0,
\label{eq:mle}
\end{equation}
with $\delta = 3[\Omega(\rho_c\theta^n) - 1] \le 0$ and
$d\delta/d\xi = (d\delta/d\theta)\,\theta'$, and the regular center
condition $\theta \simeq 1 - \xi^2/(2[3+\delta_0])$. Equation~(\ref{eq:mle})
is the same equation solved in Paper~I; only the polytropic index and the
scale relations differ.

\subsection{White dwarf model and validity range}

Cold, non-relativistic degenerate matter gives an $n = 3/2$ polytrope
with the physical constant
\begin{equation}
P = K\rho^{5/3}, \qquad
K = \frac{1.0036\times10^{13}}{\mu_e^{5/3}}~\text{(cgs)},
\end{equation}
for which we take $\mu_e = 2$ (carbon--oxygen composition)
\cite{Chandrasekhar1939,HamadaSalpeter1961}. White dwarfs form a
one-parameter family in central density: for each $\rho_c$,
$a^2 = (n{+}1)K\rho_c^{1/n-1}/4\pi G$, $R = a\,\xi_1$, and
$M = 4\pi a^3\rho_c\!\int\!\xi^2\theta^n d\xi$.

The non-relativistic approximation becomes inadequate as the electron
Fermi momentum approaches relativistic values, at
$\rho \sim 2\times10^{6}~{\rm g\,cm^{-3}}$ for $\mu_e=2$.
The retained numerical grid spans
$\rho_c \in [10^5,3\times10^6]~{\rm g\,cm^{-3}}$, with standard-model
mass labels $M_0\simeq0.16$--$0.86M_\odot$. Its upper portion therefore
extends beyond that estimate and must be read as a mathematical
polytrope comparison, not a controlled non-relativistic stellar prediction.
There is no solar-style mass or radius recalibration: $K$, $\mu_e$,
and the input central density are held fixed in each pair. Finite-temperature,
Coulomb, and envelope corrections are neglected; all results are differential
(modified vs.\ standard model at identical $\rho_c$ and EOS), which
does not establish cancellation of these inadequacies in the response. Because the $n = 3/2$
mass--density relation $M \propto \rho_c^{1/2}$ is monotonic with no
maximum, this framework makes \emph{no} statement about the Chandrasekhar
mass; a meaningful $M_{\rm Ch}$ analysis requires the relativistic
equation of state and is outside this calculation.

\subsection{Observable convention}

White dwarf tests compare radii at fixed mass. With the standard
$n = 3/2$ relation $d\ln R/d\ln M = -1/3$, the fixed-mass radius shift
can be written using the fixed-$\rho_c$ shifts as
\begin{equation}
\Delta\ln R\big|_M = \Delta\ln R\big|_{\rho_c}
+ \tfrac{1}{3}\,\Delta\ln M\big|_{\rho_c}.
\label{eq:fixedM}
\end{equation}
This identity is exact for comparison at the modified mass $M_A(\rho_c)$:
$R_0(M)\propto M^{-1/3}$ implies
$\Delta\ln R|_M=\ln[R_A(\rho_c)/R_0(M_A(\rho_c))]$.
It is not merely a first-order correction within this polytropic model.
The archived figure used $M_0(\rho_c)$ as its horizontal label; the new figure uses $M_A(\rho_c)$.
A comparison at that baseline target mass requires solving for a different
modified central density. In six audited examples this changes the log
radius shift by at most $1.6\times10^{-5}$.
The exact fractional shift is $\exp(\Delta\ln R|_M)-1$.

%============================================================================
\section{Numerical implementation and validation}
The new propagation uses the flux-form implementation audited against the
stated operator in the preceding consistency review. With
\begin{equation}
\begin{gathered}
w=\xi^{2+\delta(\theta)},\quad q=-w\theta',\\
\theta'=-q/w,\quad q'=w\theta^{3/2},\quad I'=\xi^2\theta^{3/2},
\end{gathered}
\end{equation}
we integrate ordinary mass as an ODE state. Weighted flux $q$ is not
ordinary matter mass for $A\ne0$. The full density law is retained without
suppressing its derivative near the surface. RK45 uses relative/absolute
tolerances $2\times10^{-11}/2\times10^{-13}$, maximum step .01 and
central start $\xi_0=10^{-6}$. Regular seeds are
$\theta=1-\xi_0^2/[2(3+\delta_c)]$,
$q=\xi_0^{3+\delta_c}/(3+\delta_c)$ and $I=\xi_0^3/3$.
Solver success, a terminal surface and finite outputs are required.

The null reproduces $\xi_1=3.6537537362$ and $I=2.7140551222$.
The earlier audit reproduced all 108 archived paired models using this
independently transcribed flux system. It also checked analytic mass/radius
scalings, unit changes and central-start convergence. Its strict replay
comparison with the old production archive missed a $2\times10^{-8}$
criterion (difference $6.42\times10^{-8}$), and four tightened derivative-
cutoff production solves failed. Those historical issues are retained;
the new run does not silently repair or relabel the archived data.

The new inputs are all 65 source-precision solar-a endpoints plus the local
solar-envelope peak, with amplitudes at fractions $0,.25,.5,.75,1$ of each
endpoint. Each is evaluated at 12 logarithmic central densities between
$10^5$ and $3\times10^6~{\rm g\,cm^{-3}}$. All 3960 model rows complete.
At 12 selected boundary/density pairs, tightening the flux solver by ten
and reducing the maximum step to .003 changes the exact fractional radius
shift by at most $2.78\times10^{-10}$; moving the central start to $10^{-4}$
changes it by at most $8.05\times10^{-11}$. No negative-sign requirement
is used to accept a model. These comparisons support the sampled effect
size; they do not validate the EOS approximation or prove domain extrema.

\section{Results}
%============================================================================

\subsection{Sampled conditional signatures}
\begin{figure*}[t]
\centering
\includegraphics[width=.94\textwidth]{radius_shifts.pdf}
\caption{Selected curves from the 66 nominal solar-a boundary points.
The abscissa is the actual modified comparison mass $M_A(\rho_c)$; the
ordinate is the exact percentage shift $100[\exp(\Delta\ln R|_{M_A})-1]$.
Square markers identify $\rho_c>2\times10^6~{\rm g\,cm^{-3}}$, beyond
the stated non-relativistic EOS estimate. The curves are polytropic samples,
not certified continuous-domain bounds or measured white-dwarf deficits.}
\label{fig:bands}
\end{figure*}
The largest sampled magnitude over the 3960 models is
\begin{equation}
\left.\frac{\Delta R}{R}\right|_{M_A}=-.009549333,
\end{equation}
at $\alpha=.5$, $A=.006734928$, $\rho_c=3\times10^6~{\rm g\,cm^{-3}}$
and $M_A=.84832117M_\odot$. The standard mass at that input density is
$.85925632M_\odot$, illustrating why the mass-coordinate correction matters.
The entire boundary set spans comparison masses $.15502$--$.85539M_\odot$.
The largest-deficit point is a mathematical high-density polytrope example,
not a controlled non-relativistic prediction for a real white dwarf.
\begin{table}[t]
\centering
\caption{Selected boundary samples at $\rho_c=3\times10^6~{\rm g\,cm^{-3}}$.
The last column is an exact percentage shift at $M_A$, not a log shift.
All rows in this table exceed the non-relativistic applicability estimate.}
\input{boundary_table.tex}
\label{tab:boundary}
\end{table}
At every sampled exponent and density the radius shift decreases over the
five amplitude fractions; the largest sampled magnitude occurs at a boundary
point. This is a finite-grid observation, not a proof of monotonicity between
samples or over all accepted parameters. The zero-amplitude models have zero
shift and arbitrarily small effects remain permitted. The new numbers come
from the revised solar region, rather than rescaling the older .95\% example.

The curves vary moderately with mass under this polytropic prescription.
That behavior does not follow from saturation over the full exponent range:
small-exponent models remain appreciably unsaturated at white-dwarf densities.
The full density law and its operator contribution are retained.

\subsection{Sign and physical interpretation}

The reported sampled models have radius deficits at fixed mass.
The term $(2+\delta)/\xi$ alone does not prove a universal sign because
$\delta'\ln\xi$ and calibration also contribute. Paper~I holds the
physical solar radius fixed and shows both signs of the shift in
$\xi_1$; it does not establish a physical solar contraction. V4's force
multiplier $F=(d_{\rm eff}-1)/2$ is a different prescription. A sign
change between prescriptions is not by itself a validation criterion.

\subsection{Why the V4 results were invalid}
\label{sec:v4}

For the record, the V4 implementation failed independently of the ansatz
choice: (i)~it set the polytropic constant $K = 1$ in cgs units, making
its model stars $\sim$10~m objects ($R = 430$--$926$~cm over the scanned
densities); (ii)~it began integration at $r_0 = 10^3$~cm---\emph{outside}
those stars---seeding the integration with a uniform-density sphere that
dominated the result (computed masses 9--79$\times$ the analytic values,
and mass--radius scalings inconsistent with $R \propto M^{-1/3}$);
(iii)~its ``Chandrasekhar mass shifts'' were maxima over a scan grid of a
monotonic $M(\rho_c)$ relation---the $n = 3/2$ polytrope has no maximum
mass---and so measured the grid edge, not physics. The validation gates
of this version were designed directly against these failure modes.

%============================================================================
\section{Discussion}
%============================================================================

\subsection{Observability}

The mathematical sampled scale of roughly one percent can be compared with the precision of
white dwarf mass--radius measurements. Eclipsing double-lined binaries
deliver model-independent masses and radii at the 1--2\% level
\cite{Parsons2017}, and \emph{Gaia} parallaxes constrain the
photometric mass--radius relation for large samples
\cite{Tremblay2017,Fontaine2001}. Two practical points temper the test.
First, the sampled deficit varies moderately with mass and can be
degenerate with global systematics of the mass--radius relation
(composition, envelope thickness, finite temperature) rather than
appearing as a distinctive shape; distinguishing it requires either
absolute accuracy at the few$\times10^{-3}$ level or a composition- and
temperature-controlled differential sample. Second, these reported sample scales are not lower bounds: the model
permits arbitrarily small shifts.
A quantitative exclusion or detection forecast requires composition,
thermal and envelope uncertainties, correlated measurement errors, and
checked modified stellar responses within the specified EOS and structure assumptions. A null sample does not automatically
provide a stronger bound than the solar proxy.

\subsection{Relation to the series}

The Sun supplies a conditional model/proxy threshold. White dwarfs
and neutron stars probe different structural and relativistic assumptions;
sharing a density law alone neither validates their responses nor removes
EOS degeneracy. Saturation and sign must be checked over the actual domain.

%============================================================================
\section{Conclusions}
The revised nominal solar-a sets produce sampled fixed-mass radius deficits
up to .954933\% in the $n=3/2$ model. The new figure uses $M_A$ and exact
fractional shifts. Sampling interior amplitude fractions supports a monotonic
trend on the tested grid, but does not certify a global continuum envelope.
The largest-deficit point exceeds the non-relativistic applicability estimate;
realistic EOS, thermal, composition and envelope responses remain separate tasks.

The flux-form numerical checks support the size of these conditional model
responses. V4's quantitative results remain superseded, and this polytrope
has no Chandrasekhar maximum. No detection, guaranteed observational sensitivity
or calibrated white-dwarf confidence bound is inferred from the solar proxy.

\section*{Code and revision availability}
Producing scripts, inputs, saved results and numerical audit evidence are
publicly available in the \href{https://github.com/chrismerrill1974/whitedwarfs_code}{White Dwarfs code repository}.
The accompanying source package, \nolinkurl{whitedwarfs_code}, includes the
manuscript, dependencies and resource-controlled reproduction commands.
Its \texttt{reproduction/CAPSULE.json} records file hashes;
\texttt{reproduce.py} verifies the package and rebuilds the paper from retained
scientific outputs. Full scientific runs are separate from this saved-result
reproduction. Historical scripts and results remain separately identified.

\section*{Acknowledgments}

This consistency revision was developed with computational collaboration
from Astra (Open AI) and Claude (Anthropic). The superseded V4 analysis
informed the failure-detection gates used here. We thank the developers of
\texttt{numpy}, \texttt{scipy}, and \texttt{matplotlib}.

%============================================================================
\begin{thebibliography}{9}
%============================================================================

\bibitem{PaperI} C.~K.~Merrill, Claude (Anthropic), and Astra (Open AI),
``Radial Dimensionality Theory I: Solar Neutrino and Helioseismic Constraints on Density-Dependent Spatial Dimensionality,''
revised version, November 12, 2026.

\bibitem{Chandrasekhar1939}
S.~Chandrasekhar,
\emph{An Introduction to the Study of Stellar Structure}
(University of Chicago Press, Chicago, 1939).

\bibitem{HamadaSalpeter1961}
T.~Hamada and E.~E.~Salpeter,
``Models for zero-temperature stars,''
Astrophys. J. \textbf{134}, 683 (1961).

\bibitem{Parsons2017}
S.~G.~Parsons \textit{et al.},
``Testing the white dwarf mass--radius relationship with eclipsing
binaries,''
Mon. Not. R. Astron. Soc. \textbf{470}, 4473 (2017).

\bibitem{Tremblay2017}
P.-E.~Tremblay \textit{et al.},
``The Gaia DR1 mass--radius relation for white dwarfs,''
Mon. Not. R. Astron. Soc. \textbf{465}, 2849 (2017).

\bibitem{Fontaine2001}
G.~Fontaine, P.~Brassard, and P.~Bergeron,
``The potential of white dwarf cosmochronology,''
Publ. Astron. Soc. Pac. \textbf{113}, 409 (2001).

\end{thebibliography}

\end{document}
